\section{Potenziale des Einsatzes in agilen Entwicklungsprozessen}\label{sec:potenziale}
Die berichteten Vorteile liegen auf zwei Ebenen: in der Geschwindigkeit der Erstellung und in der Qualitätssicherung samt Wissensweitergabe. Beide werden am Raster aus Kapitel~\ref{subsec:agil} geprüft.

\subsection{Produktivitäts- und Geschwindigkeitseffekte}\label{subsec:produktivitaet}
Ein kontrolliertes Experiment mit 95 rekrutierten Programmierern zeigt den Effekt unter Laborbedingungen: Bei der Implementierung eines \acs{HTTP}-Servers in JavaScript arbeitete die Gruppe mit Assistenz 55,8\,\% schneller, bei einem \gls{konfidenzintervall} von 21\,\% bis 89\,\% und dem stärksten Effekt bei weniger erfahrenen Teilnehmern.\footcite[\vglf][\pagef 2 f.]{Peng.2023} Ein Feldexperiment in einem Technologiekonzern bestätigt die Richtung unter Praxisbedingungen: Nach Einführung eines internen Sprachmodells stieg der Code-Ausstoß um 55\,\%, bei Beschäftigten mit höchstens einem Jahr Erfahrung um 67\,\%, während der Effekt bei erfahreneren Beschäftigten nicht signifikant war.\footcite[\vglf][\pagef 5]{Gambacorta.2024} Gemessen wird dort die Menge produzierter Codezeilen, nicht deren Qualität; nur 11 bis 18 Prozentpunkte des Zuwachses entfallen auf unmittelbar übernommene Modellausgaben, die übrigen 36 bis 43 führen die Autoren auf Zeitersparnis zurück.\footcite[\vglf][\pagef 5]{Gambacorta.2024} Dem steht ein randomisierter Versuch entgegen, in dem 16 erfahrene \glslink{opensource}{Open-Source}-Entwickler 246 Aufgaben in ihren eigenen \glslink{repository}{Repositorien} bearbeiteten: Sie benötigten mit Assistenz 19\,\% mehr Zeit, hatten vorab aber eine Beschleunigung um 24\,\% erwartet.\footcite[\vglf][\pagef 2]{Becker.2025} Die Befunde widersprechen einander nur scheinbar, denn sie unterscheiden sich systematisch in Erfahrung und Vertrautheit mit der Codebasis. Daraus folgt für diese Arbeit: Der Zeitgewinn konzentriert sich dort, wo Kontextwissen fehlt. Ein pauschaler Zuwachs der \gls{velocity}, also der je Iteration fertiggestellten \glspl{storypoint},\footcite[\vglf][\pagef 132]{Preussig.2020} ist nicht zu erwarten.

\subsection{Qualitätssicherungs- und Wissenseffekte}\label{subsec:qualitaetswissen}
Auf Prozessebene verbindet der DORA-Report eine um 25\,\% höhere KI-Adoption mit einer um 7,5\,\% besseren Dokumentationsqualität, einer um 3,4\,\% höheren Codequalität und um 3,1\,\% schnelleren Code Reviews.\footcite[\vglf][\pagef 37]{DORA.2024} Diese Größen lassen sich dem Raster aus Kapitel~\ref{subsec:agil} zuordnen: Die Review-Geschwindigkeit betrifft die Praktik Code Review unmittelbar, die Codequalität den Maßstab der Definition of Done. Eine Befragung in zwei großen agilen Organisationen stützt das aus Nutzersicht: Zwei Drittel berichten geringeren Aufwand für wiederkehrende Aufgaben, drei Fünftel weniger Zeit für die Informationssuche; häufiger im Arbeitsfluss bleibt jedoch nur knapp ein Drittel -- deutlich weniger als in einer früheren Herstellerbefragung.\footcite[\vglf][\pagef 46]{Wivestad.2025} In agilen Teams reichen die Anwendungen über das Programmieren hinaus: Von 21 dokumentierten Anwendungsfällen entfallen sieben auf Dokumentation und sechs auf kreative Aufgaben, darunter die Vorbereitung von \glspl{retrospektive}; zugleich ersetzt das Werkzeug die Recherche in Foren.\footcite[\vglf][\pagef 298--300]{Neumann.2026} Das senkt Einstiegshürden, bleibt aber an die einzelne Person gebunden: Dieselbe Untersuchung beschreibt die Werkzeuge als persönliche Produktivitätsassistenz, die einzelne Tätigkeiten unterstützt, ohne die Zusammenarbeit im Team zu verändern.\footcite[\vglf][\pagef 299]{Neumann.2026}
